\section{开普勒三定律与万有引力定律}\label{sec:8-4}

$17$ 世纪初，德国天文学家、数学家开普勒，在他的老师第谷（Tycho，1546---1601）对行星作的 $20$ 多年的观测的基础上，加上自己的观测与对数据的分析，总结出行星运动的三大定律．

\textbf{开普勒第一定律}\quad 行星绕太阳运行的轨道是椭圆，太阳的位置在椭圆的一个焦点上．

\textbf{开普勒第二定律}\quad 运动中，太阳至行星的向径，等时扫过等面积．

\textbf{开普勒第三定律}\quad 太阳系中各行星椭圆轨道的长轴的立方与其公转周期平方之比是一个与行星无关的常数．

行星的运动为什么会符合这些定律？几乎与开普勒同时，伽利略通过观察与分析，已得到了``惯性定律''与``自由落体定律''，牛顿也总结出运动的第二定律，问题就化成了，行星与太阳间受什么力的作用，使其运动符合开普勒三大定律？经过对开普勒、伽利略工作的深入研究，牛顿研究了月球绕地球的运动，发现月球受地球作用的力与地球上使物体下落的重力是一回事，从而提出行星受太阳作用的力也是这种力------万有引力．他于 $1687$ 年发表了震惊世界的万有引力定律．

\textbf{万有引力定律}\quad 任何两个物体，都存在一种互相吸引的力，作用在两个物体的联线上，它的大小和两个物体的质量成正比，与两物体的距离平方成反比，即
\[
F = -G\frac{m_1 m_2}{r^2},
\]
其中 $m_1$，$m_2$ 分别是两物体的质量，$r$ 为两物体间的距离，$G$ 是对任何物体都一样的常数，称为万有引力系数．

开普勒三大定律是观测结果，万有引力定律在当时的条件下无法准确测定．但假如能用数学，从开普勒三大定律推导出万有引力定律，反过来又从万有引力定律推导出开普勒三大定律，这些定律的正确性不就得到验证了么．而行星运动规律的内在根源不也就被揭示出来了么！正是牛顿在他的不朽著作《自然哲学的数学原理》中，用他刚建立起来的微积分完成了这两方面的工作，开创了微积分发展的新时代，也开创了天体力学、力学、引力理论发展的新时代．下面我们来作这两个方向的推导，当然，我们并不是重复牛顿的方式，而尽量用最简单的方法与符号．

\subsection{从开普勒定律推导万有引力定律}

取行星绕日运动的椭圆轨道所在的平面为坐标平面，太阳的位置为坐标原点，椭圆的长轴为 $x$ 轴，这时太阳至行星的向量 $\boldsymbol{r}$ 就是向径．向径有直角坐标表示与极坐标表示
\[
\boldsymbol{r} = (x,y) = (r\cos\theta, r\sin\theta),
\]
其中 $r = |\boldsymbol{r}|$ 是向径的长度．注意行星运行时，这里的 $r, x, y, \theta$ 都是 $t$ 的函数．根据极坐标下椭圆的方程，则开普勒第一定律说
\[
r = \frac{b^2}{a(1 - \varepsilon\cos\theta)},
\]
其中 $a$ 是长半轴，$b$ 是短半轴，$2c$ 为焦点间的距离：$c^2 = a^2 - b^2$，$\varepsilon = \frac{c}{a}$ 称为椭圆的偏心率，$0 < \varepsilon < 1$．

根据极坐标下平面图形的面积计算公式，行星从 $\theta_0$ 运动至 $\theta_1$ 所扫过的面积应为
\[
A = \int_{\theta_0}^{\theta_1}\frac{1}{2}r^2\,\mathrm{d}\theta.
\]
注意这里 $\theta_0$ 取定后，$\theta_1$ 随 $t$ 变化，这就是行星运动过程中向径扫过的面积的变化规律．开普勒第二定律说，等时扫过等面积，也就是说面积的变化率为常数，即
\[
\frac{\mathrm{d}A}{\mathrm{d}t} = \text{常数}.
\]
而
\[
\mathrm{d}A = \frac{1}{2}r^2\,\mathrm{d}\theta,
\]
因此
\[
\frac{\mathrm{d}A}{\mathrm{d}t} = \frac{1}{2}r^2\frac{\mathrm{d}\theta}{\mathrm{d}t} = \frac{1}{2}r^2\dot{\theta} = \frac{1}{2}r^2\omega = \text{常数}.
\]
这里 $\omega = \frac{\mathrm{d}\theta}{\mathrm{d}t}$ 是向径的角速度．上式实际上就是说，行星在运动中角动量守恒（为常数）．

让我们回忆什么叫角动量．牛顿定律可以写成
\[
\boldsymbol{f} = m\boldsymbol{a} = \frac{\mathrm{d}}{\mathrm{d}t}(m\boldsymbol{v}),
\]
$m\boldsymbol{v}$ 称为物体的动量．因此牛顿定律可以解释为物体受的力等于动量的变化率．在等式两边用向量 $\boldsymbol{r}$ 作叉乘（解析几何中的向量积），得
\[
\boldsymbol{r}\times\boldsymbol{f} = m(\boldsymbol{r}\times\boldsymbol{a}) = m\left(\boldsymbol{r}\times\frac{\mathrm{d}\boldsymbol{v}}{\mathrm{d}t}\right) = \frac{\mathrm{d}}{\mathrm{d}t}(m\boldsymbol{r}\times\boldsymbol{v}).
\]
（这里用到了叉乘的微商法则
\[
\frac{\mathrm{d}}{\mathrm{d}t}(\boldsymbol{r}\times\boldsymbol{v}) = \frac{\mathrm{d}\boldsymbol{r}}{\mathrm{d}t}\times\boldsymbol{v} + \boldsymbol{r}\times\frac{\mathrm{d}\boldsymbol{v}}{\mathrm{d}t} = \boldsymbol{v}\times\boldsymbol{v} + \boldsymbol{r}\times\frac{\mathrm{d}\boldsymbol{v}}{\mathrm{d}t} = \boldsymbol{r}\times\frac{\mathrm{d}\boldsymbol{v}}{\mathrm{d}t},
\]
它是很容易证明的，另外还用到 $\boldsymbol{v}\times\boldsymbol{v} = \boldsymbol{0}$．我们看到，等式左边是力矩，对比力等于动量的变化率，人们自然称 $m\boldsymbol{r}\times\boldsymbol{v}$ 为角动量．则上式说，物体运动所受的力矩等于角动量的变化率．）

下面我们来计算行星运动的角动量，这时
\[
\boldsymbol{r} = (r\cos\theta, r\sin\theta, 0),
\]
则
\[
\boldsymbol{v} = \frac{\mathrm{d}\boldsymbol{r}}{\mathrm{d}t} = (\dot{r}\cos\theta - r\sin\theta\cdot\dot{\theta},\, \dot{r}\sin\theta + r\cos\theta\cdot\dot{\theta}, 0).
\]
用解析几何的叉乘公式
\[
m\boldsymbol{r}\times\boldsymbol{v} = m\begin{vmatrix}
\boldsymbol{i} & \boldsymbol{j} & \boldsymbol{k} \\
r\cos\theta & r\sin\theta & 0 \\
\dot{r}\cos\theta - r\sin\theta\cdot\dot{\theta} & \dot{r}\sin\theta + r\cos\theta\cdot\dot{\theta} & 0
\end{vmatrix} = m r^2\dot{\theta}\boldsymbol{k}.
\]
开普勒第二定律说，$r^2\dot{\theta}$ 是常数，也就是说角动量守恒．而力矩是角动量的变化率，也就是行星运动过程中，力矩为 $0$，即所受的力 $\boldsymbol{f}$ 与 $\boldsymbol{r}$ 共线．这就证明了行星所受的力在行星与太阳的联线上．

下面来求这力的大小，为此把加速度向量分解为平行向径与垂直向径的两部分：
\[
\boldsymbol{a} = \boldsymbol{a}_r + \boldsymbol{a}_\theta,
\]
其中 $\boldsymbol{a}_r$ 方向的单位向量为 $\boldsymbol{e}_r = (\cos\theta, \sin\theta)$，$\boldsymbol{a}_\theta$ 方向的单位向量为 $\boldsymbol{e}_\theta = (-\sin\theta, \cos\theta)$．在前面 $\boldsymbol{v}$ 的表达式中继续对 $t$ 求一次微商，得
\begin{align*}
\boldsymbol{a} = \frac{\mathrm{d}\boldsymbol{v}}{\mathrm{d}t} &= \big(\ddot{r}\cos\theta - 2\dot{r}\sin\theta\cdot\dot{\theta} - r\cos\theta\cdot\dot{\theta}^2 - r\sin\theta\cdot\ddot{\theta}, \\
&\qquad \ddot{r}\sin\theta + 2\dot{r}\cos\theta\cdot\dot{\theta} - r\sin\theta\cdot\dot{\theta}^2 + r\cos\theta\cdot\ddot{\theta}\big) \\
&= (\ddot{r} - r\omega^2)(\cos\theta, \sin\theta) + (2\dot{r}\dot{\theta} + r\ddot{\theta})(-\sin\theta, \cos\theta) \\
&= (\ddot{r} - r\omega^2)\boldsymbol{e}_r + (2\dot{r}\dot{\theta} + r\ddot{\theta})\boldsymbol{e}_\theta.
\end{align*}
由 $r^2\dot{\theta} = \text{常数}$，知
\[
2r\dot{r}\dot{\theta} + r^2\ddot{\theta} = r(2\dot{r}\dot{\theta} + r\ddot{\theta}) = 0,
\]
从而知 $a_\theta = 0$，故 $\boldsymbol{a}$ 的大小等于 $\ddot{r} - r\omega^2$．由开普勒第二定律知
\[
\pi a b = \int_0^T \frac{1}{2}r^2\omega\,\mathrm{d}t = \frac{1}{2}r^2\omega T,
\]
其中 $T$ 是行星运动的周期，即绕日旋转一周所需的时间，因此
\begin{align}
r^2\omega &= \frac{2\pi a b}{T}, \label{eq:8-4-1}\\
r^2\omega^2 &= \left(\frac{2\pi a b}{r T}\right)^2. \notag
\end{align}
用 $r = \frac{b^2}{a(1-\varepsilon\cos\theta)}$ 去除，得
\[
r\omega^2 = \frac{4\pi^2 a^3}{T^2}\frac{1}{r^2}(1-\varepsilon\cos\theta).
\]
另一方面由
\[
\frac{1}{r} = \frac{a}{b^2}(1-\varepsilon\cos\theta),
\]
有
\[
-\frac{\dot{r}}{r^2} = \frac{a}{b^2}\varepsilon\sin\theta\cdot\omega.
\]
因此由 \eqref{eq:8-4-1} 得
\[
-\dot{r} = \frac{a}{b^2}\varepsilon\sin\theta\cdot r^2\omega = \frac{2\pi a b}{T}\frac{a}{b^2}\varepsilon\sin\theta.
\]
再用 \eqref{eq:8-4-1}，有
\[
-\ddot{r} = \frac{2\pi a b}{T}\frac{a}{b^2}\varepsilon\cos\theta\cdot\omega = \frac{4\pi^2 a^3}{T^2}\frac{\varepsilon\cos\theta}{r^2}.
\]
结果得到
\[
\boldsymbol{f} = \boldsymbol{f}_r = m(\ddot{r} - r\omega^2)\boldsymbol{e}_r = -\frac{4\pi^2 a^3}{T^2}\frac{m}{r^2}\boldsymbol{e}_r.
\]
由开普勒第三定律，$\frac{a^3}{T^2}$ 为常数，令
\[
G = \frac{4\pi^2 a^3}{T^2 M},
\]
其中 $M$ 为太阳质量，则 $G$ 是与行星无关由太阳唯一决定的常数，从而行星受的力为
\[
\boldsymbol{f} = -G\frac{M m}{r^2}\boldsymbol{e}_r.
\]
这就推导出了万有引力．

\subsection{从万有引力推导开普勒定律}

设行星受的力为万有引力，即
\[
\boldsymbol{f} = -G m M \frac{1}{r^3}(x,y,z),
\]
由
\[
\boldsymbol{0} = \boldsymbol{r}\times\boldsymbol{f} = \frac{\mathrm{d}}{\mathrm{d}t}(\boldsymbol{r}\times m\boldsymbol{v}),
\]
故
\[
\boldsymbol{r}\times\boldsymbol{v} = \text{常向量}.
\]
根据前面的计算 $|\boldsymbol{r}\times\boldsymbol{v}| = r^2\omega = r^2\frac{\mathrm{d}\theta}{\mathrm{d}t}$，把它和极坐标面积计算公式联系起来，这就是开普勒第二定律：向径等时扫过等面积．

由 $\boldsymbol{r}\times\boldsymbol{v} = \text{常向量}$，可取此向量为 $z$ 轴，则行星在一与此向量垂直的平面上运动，这时
\[
\boldsymbol{f} = -G m M \frac{1}{r^3}(x,y),
\]
其中 $\boldsymbol{r} = (x,y)$．由于
\[
\frac{\mathrm{d}}{\mathrm{d}t}\left(\frac{1}{2}m|\boldsymbol{v}|^2\right) = \frac{\mathrm{d}}{\mathrm{d}t}\left(\frac{1}{2}m\boldsymbol{v}\cdot\boldsymbol{v}\right) = m\boldsymbol{v}\cdot\boldsymbol{a} = \boldsymbol{v}\cdot\boldsymbol{f},
\]
取积分得
\[
\frac{1}{2}m v_1^2 - \frac{1}{2}m v_0^2 = \int_{t_0}^t \boldsymbol{v}\cdot\boldsymbol{f}\,\mathrm{d}t = -\int_{t_0}^t G m M \frac{x\dot{x} + y\dot{y}}{r^3}\,\mathrm{d}t = U_1 - U_0,
\]
其中
\[
U = \frac{G m M}{r}
\]
称为位能，而 $\frac{1}{2}m v^2$ 称为动能．上述等式说明，在向心力作用下，行星的动能加位能恒等于常数：
\[
\frac{1}{2}m(\dot{x}^2 + \dot{y}^2) - \frac{G m M}{r} = C.
\]
注意到
\[
x = r\cos\theta,\quad y = r\sin\theta,\quad \dot{x}^2 + \dot{y}^2 = \dot{r}^2 + r^2\dot{\theta}^2,
\]
即把运动方程化为
\[
\begin{cases}
\dfrac{1}{2}m\left(\dot{r}^2 + (r\omega)^2\right) - \dfrac{G m M}{r} = C,\\[8pt]
r^2\omega = h,
\end{cases}
\]
其中 $C, h$ 是常量．设 $\omega = \dot{\theta} \neq 0$（否则 $\theta = \text{常数}$，行星作直线运动），则由
\[
\omega = \frac{h}{r^2},\qquad \dot{r} = \frac{\mathrm{d}r}{\mathrm{d}\theta}\omega,
\]
知
\[
\frac{m}{2}\left[\frac{h^2}{r^4}\left(\frac{\mathrm{d}r}{\mathrm{d}\theta}\right)^2 + \frac{h^2}{r^2}\right] - \frac{G m M}{r} = C,
\]
或
\[
\left(\frac{\mathrm{d}r}{\mathrm{d}\theta}\right)^2 = r^4\left[\frac{2C}{m h^2} + \frac{2G M}{h^2}\frac{1}{r} - \frac{1}{r^2}\right].
\]
令
\[
r = \frac{1}{W},\quad \frac{1}{p} = \frac{G M}{h^2},\quad \varepsilon^2 = 1 + \frac{2C h^2}{m G^2 M^2},
\]
则
\[
\left(\frac{\mathrm{d}W}{\mathrm{d}\theta}\right)^2 = \frac{\varepsilon^2}{p^2} - \left(W - \frac{1}{p}\right)^2.
\]
令 $V = W - \frac{1}{p}$，积分便得
\[
\theta - \theta_0 = \int \frac{\mathrm{d}V}{\sqrt{\frac{\varepsilon^2}{p^2} - V^2}},
\]
从而
\[
V = \frac{\varepsilon}{p}\sin(\theta - \theta_0),
\]
或
\[
\frac{1}{r} - \frac{1}{p} = \frac{\varepsilon}{p}\sin(\theta - \theta_0).
\]
取 $\theta_0 = \frac{\pi}{2}$，便得
\[
r = \frac{p}{1 - \varepsilon\cos\theta},
\]
其中
\[
\begin{cases}
p = \dfrac{b^2}{a} = \dfrac{h^2}{G M},\\[8pt]
\varepsilon^2 = 1 + \dfrac{2C h^2}{m G^2 M^2}.
\end{cases}
\]
不难看出，当 $C < 0$ 时，$\varepsilon < 1$，轨道为椭圆，这就是开普勒第一定律．当 $C = 0$ 时，$\varepsilon = 1$，轨道是抛物线．当 $C > 0$ 时，$\varepsilon > 1$，轨道是双曲线．后两种情形表明，行星受太阳的万有引力作用，运动轨道也可能是抛物线与双曲线．这一点是开普勒没有观测到而以后被一些小行星与彗星的运行轨道证实了的．

由 $2\pi a b = \int_0^T r^2\omega\,\mathrm{d}t = hT$，知
\[
4\pi^2 a^2 b^2 = h^2 T^2,
\]
而
\[
\frac{b^2}{a} = \frac{h^2}{G M},
\]
故
\[
\frac{T^2}{a^3} = \frac{4\pi^2}{G M},
\]
右边是与行星无关的常数，这就推出了开普勒第三定律．
